\chapter{Selbstgerechtigkeit}
\quellseite{9}

Eine alte, schwer leidende Frau sagte zu ihrem sie besuchenden Seelsorger: „Ich finde bei mir keine so große Sünde, womit ich diese Züchtigung verdient haben könnte. Es muß die Sünde meiner Eltern oder Großeltern sein, für welche ich heimgesucht werde.“ — So groß ist der selbstsüchtige Stolz des selbstgerechten Herzens, daß man lieber noch einen Tadel wirft auf Eltern und Vorfahren, als daß man die eigene Schuld zugesteht. Oefter noch muß Gott die Schuld tragen, weil er uns so schwach geschaffen und weil er uns in so große Versuchungen habe kommen lassen. Adam wälzt die Schuld — die eheliche
\quellseite{10}
Treue verletzend — auf Eva und zuletzt — Gottes Heiligkeit und Güte lästernd — auf Gott, der ihm Eva zugesellt habe, und Eva schiebt sie auf die Schlange, welche doch Gott ins Paradies gesetzt habe, 1 Mos. 3, 12. f. (Sonntagsblatt fürs Haus.)
